\documentclass[12pt,a4paper]{article}
\usepackage[margin=2.5cm]{geometry}
\usepackage{amsmath,amssymb}
\usepackage{booktabs}
\usepackage{hyperref}
\usepackage{xcolor}
\usepackage{caption}
\usepackage{float}
\usepackage{placeins}
\usepackage{enumitem}
\usepackage{parskip}
\usepackage{titlesec}
\usepackage{siunitx}

% ── Colour macros ─────────────────────────────────────────────────────────────
\definecolor{ProvenGreen}{HTML}{22c55e}
\definecolor{DangerRed}{HTML}{ef4444}
\definecolor{CbfAmber}{HTML}{f59e0b}
\newcommand{\proven}[1]{\textcolor{ProvenGreen}{\textbf{#1}}}
\newcommand{\danger}[1]{\textcolor{DangerRed}{\textbf{#1}}}
\newcommand{\cbf}[1]{\textcolor{CbfAmber}{\textbf{#1}}}

% ── Section formatting ────────────────────────────────────────────────────────
\titleformat{\section}{\large\bfseries}{{\thesection}}{1em}{}[\titlerule]
\titleformat{\subsection}{\normalsize\bfseries}{{\thesubsection}}{1em}{}

% ─────────────────────────────────────────────────────────────────────────────
\title{\textbf{FleetSafe-VLA/VLN}\\[6pt]
       \large Year 1 Annual Review\\[4pt]
       \normalsize School of Computing, Newcastle University}
\author{Frank van Laarhoven\\
        \small \texttt{F.Van-Laarhoven2@newcastle.ac.uk}\quad
        ORCID: 0009-0006-8931-0364}
\date{May 2026}
% ─────────────────────────────────────────────────────────────────────────────

\begin{document}
\maketitle
\tableofcontents
\newpage

% ─────────────────────────────────────────────────────────────────────────────
\begin{abstract}
FleetSafe-VLA/VLN investigates architecture-agnostic command-layer safety
filtering under tested conditions for heterogeneous embodied visual navigation
models (GNM, ViNT, NoMaD) in hospital-like environments, without modifying the
underlying policies.

Year~1 established a reproducible evaluation infrastructure and produced a key
scientific finding: vulnerability to invisible hazards is
\emph{navigation-paradigm-dependent}. Goal-directed policies (GNM, ViNT)
collide consistently, while diffusion-based exploration (NoMaD) naturally
avoids them. FleetSafe achieves zero collisions across all three paradigms
while demonstrating selective intervention and do-no-harm behaviour, with
intervention rate equal to 0\% for NoMaD in the tested Isaac Sim scenario.

The programme is structured around five research questions and a staged
evidence taxonomy (\textsc{sim-mock} $\to$ \textsc{sim-mujoco} $\to$
\textsc{sim-isaac} $\to$ \textsc{real-proven}). Year~2 prioritises
\textsc{real-proven} physical validation and submission to CoRL~2026 or RA-L.

\medskip
\noindent\textbf{Keywords:} Control Barrier Functions, visual navigation,
architecture-agnostic safety, healthcare robotics, paradigm-dependent
robustness
\end{abstract}

% ─────────────────────────────────────────────────────────────────────────────
\section{Research Framing}
\label{sec:framing}

Embodied visual navigation models can navigate toward goals but do not
automatically provide deployment-time safety guarantees. This is critical in
hospital-style environments, where narrow corridors, dynamic human traffic, and
low tolerance for unsafe behaviour make safety a first-order requirement.

\paragraph{Why these three backbones?}
GNM represents classical CNN-based generalist navigation, ViNT is a
Transformer-based foundation model, and NoMaD introduces diffusion-based
goal-masked policies. Together they span the major contemporary paradigms in
visual navigation --- convolutional, transformer, and generative diffusion ---
making the architecture-agnostic claim both broad and falsifiable.

The central hypothesis is that a command-layer CBF-QP safety filter can enforce
collision-free and socially aware navigation across these three qualitatively
distinct architectures while preserving task-level performance and exposing
paradigm-dependent intervention behaviour.

% ─────────────────────────────────────────────────────────────────────────────
\section{Related Work and Gap}
\label{sec:related}

Recent mobile-robot foundation models such as GNM~\cite{shah2023gnm},
ViNT~\cite{majumdar2023vint}, and NoMaD~\cite{sridhar2024nomad} provide strong
nominal navigation capability. FleetSafe addresses the deployment-time safety
gap by evaluating architecture-agnostic command-layer safety filtering across
these three distinct paradigms. Prior safety approaches either require policy
modification (safe RL, shielding) or gradient access (constrained policy
optimisation); command-layer CBF filtering requires neither, making it
uniquely suited to frozen pretrained policies.

% ─────────────────────────────────────────────────────────────────────────────
\section{Theoretical Framing}
\label{sec:theory}

FleetSafe studies whether a reactive safety filter can render nominally unsafe
commands admissible with respect to a formally defined safe set while remaining
agnostic to the internal policy architecture.

Let the robot state be $x \in \mathbb{R}^n$, admissible control
$u \in \mathcal{U}$, and safety function $h(x)$. The safe set is defined as
\[
  \mathcal{C} = \{x \in \mathbb{R}^n : h(x) \geq 0\}.
\]
A Control Barrier Function requires
\[
  \dot{h}(x,u) + \alpha(h(x)) \geq 0,
\]
where $\alpha(\cdot)$ is an extended class-$\mathcal{K}$ function
\cite{ames2019cbf}. This is realised via the Quadratic Program
\[
  \min_{u} \|u - u_{\text{nom}}\|^2
  \quad \text{s.t.}\quad
  \dot{h}(x,u) + \alpha(h(x)) \geq 0,\quad u \in \mathcal{U}.
\]
For cylindrical obstacles (radius $r_i$, centre $\mathbf{o}_i$), the
surface-distance barrier function is
\[
  h_{\text{CBF}}(\mathbf{p},i)
  = \bigl(\|\mathbf{p}-\mathbf{o}_i\|_2 - r_i\bigr)^2 - d_{\text{safe}}^2,
\]
which has a closed-form single-constraint QP solution, enabling real-time
control at \SI{10}{\hertz}.

% ─────────────────────────────────────────────────────────────────────────────
\section{Methodology and Evidence Ladder}
\label{sec:methodology}

The programme follows a simulation-first, staged-validation methodology
justified by the safety-critical nature of embodied AI. Table~\ref{tab:tiers}
defines the evidence tiers. Claims may only be made at or below the tier at
which evidence exists; an automated PROVEN gate prevents premature claims.

\begin{table}[t]
\centering
\caption{Staged evidence taxonomy. Each tier unlocks progressively stronger
         publication claims.}
\label{tab:tiers}
\begin{tabular}{@{}lll@{}}
\toprule
Tier & Meaning & Purpose \\
\midrule
\textsc{sim-mock}    & Simplified backend    & Infrastructure validation, early
                                               hypothesis testing \\
\textsc{sim-mujoco}  & Physics-grounded      & Publication-grade comparative
                                               evaluation \\
\textsc{sim-isaac}   & Photoreal sim-to-sim  & Visual transfer and robustness
                                               stress testing \\
\textsc{real-manual} & Manual commissioning  & Motion confirmation and safe
                                               deployment preparation \\
\textsc{real-proven} & Scripted physical     & Repeatable real-world validation \\
\bottomrule
\end{tabular}
\end{table}

% ─────────────────────────────────────────────────────────────────────────────
\section{Year 1 Progress and Current Evidence}
\label{sec:results}

FleetSafe achieves zero collisions across all paradigms while remaining
selectively active. Goal-directed models (GNM, ViNT) require frequent
intervention; diffusion-based NoMaD demonstrates clear do-no-harm behaviour
with intervention rate equal to 0\% in the tested Isaac Sim scenario.

\subsection{Navigation Paradigms Evaluated}

\begin{table}[t]
\centering
\caption{The three navigation paradigms evaluated in FleetSafe. Together they
         span the dominant directions in current visual navigation research.}
\label{tab:paradigms}
\begin{tabular}{@{}llll@{}}
\toprule
Model & Paradigm & Core Mechanism & Key Characteristic \\
\midrule
GNM   & CNN-based generalist      & Goal-conditioned imitation   &
        Efficient cross-embodiment baseline \\
ViNT  & Transformer foundation    & Attention over visual history &
        Scalable high-capacity reasoning \\
NoMaD & Diffusion policy          & Goal-masked trajectory generation &
        Multimodal exploration and navigation \\
\bottomrule
\end{tabular}
\end{table}

\FloatBarrier
\subsection{Invisible Hazard Experiment — Isaac Sim (PROVEN)}

Table~\ref{tab:isaac} presents the primary scientific result: vulnerability to
invisible map-registered hazards is navigation-paradigm-dependent.

\begin{table}[t]
\centering
\caption{Isaac Sim hospital corridor: paradigm-dependent failure and FleetSafe
         selective intervention ($n = 50$ seeds, PROVEN gate certified).
         Collision Rate (CR) and CBF Intervention Rate (IR) reported in
         percentage. Min-dist is mean closest approach to hazard surface
         ($r = 1.0$\,m cylinder).}
\label{tab:isaac}
\setlength{\tabcolsep}{6pt}
\begin{tabular}{@{}lccccr@{}}
\toprule
Model & Paradigm & Baseline CR & FleetSafe CR & Intervention Rate &
Min.\ Dist.\ (m) \\
\midrule
GNM   & Goal-directed        & \danger{100\%} & \proven{0\%} & \cbf{49.7\%} &
        0.075 \\
ViNT  & Goal-directed        & \danger{100\%} & \proven{0\%} & \cbf{37.9\%} &
        0.089 \\
NoMaD & Diffusion/exploration & \proven{0\%}  & \proven{0\%} & \textbf{0.0\%} &
        1.531 \\
\bottomrule
\end{tabular}
\end{table}

\FloatBarrier
\subsection{System Architecture}

Figure~\ref{fig:arch} illustrates the command-layer safety stack that is
applied without modification to all three backbone policies.

\begin{figure}[htbp]
\centering
\small
\begin{verbatim}
  RGB Camera + State Estimate
              |
   +----------+----------+------------------+
   |          |          |                  |
  GNM        ViNT      NoMaD              ...
  (CNN)  (Transformer) (Diffusion)
   |          |          |
   +----------+----------+
              |
     Multi-Backbone Adapter
              |
       CBF-QP Safety Filter
       h(p) = d_surf^2 - d_safe^2
              |
       cmd_vel_safe
              |
       Robot Execution
\end{verbatim}
\caption{FleetSafe-VLA/VLN system architecture. GNM, ViNT, and NoMaD feed into
         a shared command-layer CBF-QP safety stack. The filter operates at the
         velocity-command interface and requires no access to policy internals.}
\label{fig:arch}
\end{figure}

% ─────────────────────────────────────────────────────────────────────────────
\section{Interpretation and Scientific Contribution}
\label{sec:contribution}

The results demonstrate that invisible-hazard vulnerability is
\emph{navigation-paradigm-dependent} rather than universal. Goal-directed
policies (GNM, ViNT) commit to obstacle-crossing trajectories toward their
targets and collide at 100\% RAW collision rate, while diffusion-based NoMaD
avoids the same hazard naturally (0\% RAW, min clearance $= 1.5$\,m) because
its exploration prior does not commit to a fixed forward path.

FleetSafe provides selective, non-intrusive command-layer safety filtering:
the same filter that prevents 100\% collision in GNM/ViNT corridor episodes
adds zero intervention cost to the already-safe NoMaD. This directly addresses
the over-conservatism criticism of safety filters
--- a standard concern in deployment-time safety systems~\cite{brunke2022safe}.

The central claim is strengthened by evaluating across three fundamentally
different architectures: the architecture-agnostic assertion is substantially
more meaningful when the three tested models span convolutional, transformer,
and generative-diffusion paradigms rather than minor variants of a single class.

% ─────────────────────────────────────────────────────────────────────────────
\section{Data Governance, Reproducibility, and Ethics}
\label{sec:governance}

Active data are stored on Newcastle University RDS with role-based access,
encryption, and audit logging. No patient or identifiable human data are
collected. Real-robot experiments remain conservative, including wheels-lifted
validation, operator override, and emergency-stop procedures. Future
occupied-corridor trials require full ethics and NHS governance approval.

The benchmark and code will be released publicly under MIT and CC-BY~4.0
licensing on GitHub and Hugging Face in Q2~2027. All simulation results carry
$n = 50$ seeds and automated PROVEN-gate certification; evidence artefacts
are hash-committed to the repository.

% ─────────────────────────────────────────────────────────────────────────────
\section{Year 2 Priorities and Trajectory}
\label{sec:year2}

\begin{enumerate}[leftmargin=*, label=\arabic*.]
  \item Capture scripted \textsc{real-proven} rosbag on the Yahboom M3Pro
    platform (ROS2), validating the CBF filter under real perception noise
    and actuator delay.
  \item Complete extended statistical analysis (Fisher exact tests,
    Wilson confidence intervals across all scenes and backends).
  \item Submit \textit{Paper 1} on architecture-agnostic command-layer safety
    and paradigm-dependent robustness to CoRL~2026 or RA-L.
  \item Complete the Isaac Sim photoreal stress test and open benchmark
    release on GitHub/Hugging Face.
  \item Extend evaluation to non-cylindrical obstacles (door frames,
    furniture) and dynamic agents (constant-velocity pedestrians).
\end{enumerate}

Year~1 has established the theoretical, experimental, and infrastructural
foundations for a rigorous investigation into architecture-agnostic
command-layer safety filtering for embodied visual navigation. The programme
is now positioned to deliver publication-quality contributions during Years~2
and~3.

% ─────────────────────────────────────────────────────────────────────────────
\clearpage
\begin{thebibliography}{99}

\bibitem{shah2023gnm}
D.~Shah, B.~Eysenbach, G.~Kahn, N.~Rhinehart, S.~Levine.
\textit{GNM: A General Navigation Model to Drive Any Robot}.
IEEE International Conference on Robotics and Automation (ICRA), 2023.

\bibitem{majumdar2023vint}
A.~Majumdar, K.~Yadav, S.~Arnaud, Y.~Ma, C.~Chen, S.~Silwal,
A.~Jain, V.~Berges, T.~Sukhatme, P.~Abbeel, D.~Shah.
\textit{ViNT: A Foundation Model for Visual Navigation}.
Conference on Robot Learning (CoRL), 2023.

\bibitem{sridhar2024nomad}
A.~Sridhar, D.~Shah, C.~Glossop, S.~Levine.
\textit{NoMaD: Goal Masked Diffusion Policies for Navigation and Exploration}.
IEEE International Conference on Robotics and Automation (ICRA), 2024.

\bibitem{ames2019cbf}
A.~D.~Ames, S.~Coogan, M.~Egerstedt, G.~Notomista, K.~Sreenath, P.~Tabuada.
\textit{Control barrier functions: Theory and applications}.
European Control Conference (ECC), 2019.

\bibitem{brunke2022safe}
L.~Brunke, M.~Greeff, A.~W.~Hall, Z.~Yuan, S.~Zhou, J.~Panerati,
A.~P.~Schoellig.
\textit{Safe Learning in Robotics: From Learning-Based Control to Safe
Reinforcement Learning}.
Annual Review of Control, Robotics, and Autonomous Systems,
5:411--444, 2022.

\end{thebibliography}

\end{document}
